\section*{第 \ref{ch:g9:uniform-varied-motion} 章 --- 运动：匀速与变速}

\begin{solution}{exo:g9:uniform-varied-motion:1}
在相等的时间间隔处标出物体的位置。两种
方法：逐帧查看视频，或用打点计时器在被拉动的纸带上打点
（或频闪照片）。
\end{solution}

\begin{solution}{exo:g9:uniform-varied-motion:2}
间隔相等：匀速运动。$v = 0.04 \div 0.1 = \qty{0.4}{m/s}$。
\end{solution}

\begin{solution}{exo:g9:uniform-varied-motion:3}
间隔逐渐增大：加速运动。第一段：$0.01 \div 0.1 =
\qty{0.1}{m/s}$；最后一段：$0.07 \div 0.1 = \qty{0.7}{m/s}$。
\end{solution}

\begin{solution}{exo:g9:uniform-varied-motion:4}
雷达测量的是瞬时速度 --- 也就是经过闸口那一时刻的速度，
“门到门二十分钟”算出的是平均
速度 --- 总路程除以总时间，停停走走、快跑慢跑都算在内。
\end{solution}

\begin{solution}{exo:g9:uniform-varied-motion:5}
匀速运动中，速度始终不变，因此每一段
--- 无论长短，算出的结果都相同，两个概念
合二为一。
\end{solution}

\begin{solution}{exo:g9:uniform-varied-motion:6}
每秒速度约增加 \qty{9.8}{m/s} --- 在数值上等于当地的
$g$，即 \qty{9.8}{N/kg}：本学年最巧妙的呼应。
\end{solution}

\begin{solution}{exo:g9:uniform-varied-motion:7}
平衡（或不受）的力不改变运动状态：静止的
仍然静止，匀速直线运动继续进行。不受力的
物体将永远滑行下去；客机的四方平衡 --- 推力与阻力，升力与
阻力与推力、升力与重力相互抵消 --- 这正是它
能以 \qty{900}{km/h} 平稳飞行的缘故。
\end{solution}

\begin{solution}{exo:g9:uniform-varied-motion:8}
探测器在真空中不受力，已滑行了数十年，
发动机早已冷却 --- 运动的延续无需外力。
那辆“证明”相反结论的自行车，在摩擦力
和空气阻力下减速 --- 实为伪装的不平衡反向力；在冰面上，
伪装变淡，便能滑行得很远。
\end{solution}

\begin{solution}{exo:g9:uniform-varied-motion:9}
间隔逐渐缩小：减速。速度依次为 $0.6$、$0.5$、$0.4$、
$0.3$、\qty{0.2}{m/s} --- 每个时间间隔减少 \qty{0.1}{m/s}。
若规律保持，再过两个时间段（间隔为 \qty{1.0}{cm}，随后
便为零），纸带即静止。
\end{solution}

\begin{solution}{exo:g9:uniform-varied-motion:10}
每 \qty{0.1}{s} 增加 \qty{0.1}{m/s}，即每秒增加 \qty{1}{m/s}，
约为自由下落每秒增加的 \qty{9.8}{m/s} 的十分之一，
与地球的引力相比，是一股温和的推力。
\end{solution}

\begin{solution}{exo:g9:uniform-varied-motion:11}
第一阶段：重力大于刚开始较小的空气阻力 ---
合力向下，速度增大。伞张开后：
迅速增大的空气阻力逐渐上升，直至恰好与重力平衡
--- 二力平衡，下落便稳定为
题中所说的 \qty{5}{m/s} 匀速下落。
\end{solution}

\begin{solution}{exo:g9:uniform-varied-motion:12}
记录：拍摄释放过程，每隔 \qty{0.1}{s} 取一帧（或
将小车连接到打点计时器上）；沿斜面下滑并穿过地毯，逐段测量
间隔距离。数据：各段速度，依次列出。
预期结论：第一幕，间隔逐渐增大 --- 在斜面上做加速运动
（重力沿斜面方向的不平衡分力）；第二幕，
间隔逐渐缩小 --- 在地毯上做减速运动（摩擦力的
不平衡反向力）。一个玩具，展示了这句平静陈述
的两个侧面。
\end{solution}

\begin{solution}{pb:g9:uniform-varied-motion:1}
\textbf{1.} 匀速；$v = 0.08 \div 0.1 = \qty{0.8}{m/s}
\approx \qty{2.9}{km/h}$ --- 相当于散步。
\textbf{2.} 加速运动：从 $0.02 \div 0.1 = \qty{0.2}{m/s}$
增至 \qty{1.0}{m/s}；每个时间间隔增加 \qty{0.2}{m/s} ---
即每秒增加 \qty{2}{m/s}。
\textbf{3.} 急剧减速。沙子提供了很大的
不平衡反向力 --- 即很大的摩擦力 --- 速度
急剧下降：$1.0$、$0.7$、$0.45$、$0.25$、
\qty{0.1}{m/s}。
\textbf{4.} 前三个时间间隔以 \qty{0.5}{m/s} 匀速滑行，
随后加速 --- 间隔骤增至 $7.5$ 和
\qty{10.0}{cm}：第三与第四个时间间隔之间，有什么东西开始
施力（或是斜坡，或是一推）。这一事件发生在第三
段间隔的末端。
\textbf{5.} 两种说法都对：平均速度综合了整个行程中
的等待与行驶；雷达读取的是经过时的
瞬时速度。交通法规以雷达为准 --- 限速约束的是
瞬时速度，再平缓的平均速度也不能为超速
的那一瞬间开脱。
\textbf{6.} 匀加速。相等的时间内速度增量相等，
说明有一个稳定不变的合力 --- 原因恒定，
速度增长也恒定。
\textbf{7.} 地面对鞋子的向前推力
与空气和地面对身体的向后阻力相抵，
二力平衡，因此速度稳定为 \qty{0.8}{m/s}。
\textbf{8.} 增量：每 \qty{0.1}{s} 的时间间隔增加 \qty{0.98}{m/s}
--- 即每秒 \qty{9.8}{m/s}，正是当地的 $g$，
由石块下落的纸带再现。
\textbf{9.} 该说法成立：速度依次为 $1.2$、$1.15$、$1.10$、
$1.05$、\qty{1.0}{m/s} --- 一路平缓减速。
\textbf{10.} 出口：$1.2 \times 3.6 = \qty{4.3}{km/h}$；入口
\qty{3.6}{km/h} --- 均远低于骑车人的标准：
并无超速鲁莽之嫌；操控方面，或许另当别论。
\textbf{11.} “在平地上，除向后的不平衡摩擦力外，没有什么能使滚动的滑板停下 ---
--- 光滑场地提供的摩擦力很小，
因此滑板保持了大部分运动，恰如
平衡力定律所说的那样。”
\textbf{12.} 例如：“间隔相等，速度平稳 ---
没有合力扰乱它的平稳。间隔变化则需要一个
原因：某种不平衡的推力，正如这句平静的陈述所言。”
\end{solution}
